%% Supplementary Material
%% Empirical Huff Curves for Brazil — Journal of Hydrology
\documentclass[11pt]{article}

\usepackage[a4paper,margin=2.4cm,top=2.8cm,bottom=2.8cm]{geometry}
\usepackage[authoryear]{natbib}
\usepackage{amsmath}
\usepackage{graphicx}
\usepackage{booktabs}
\usepackage{siunitx}
\usepackage{xcolor}
\definecolor{linkblue}{HTML}{1F6FB5}
\usepackage{microtype}
\usepackage{setspace}
\usepackage{fancyhdr}
\usepackage[
  colorlinks = true,
  linkcolor  = linkblue,
  citecolor  = linkblue,
  urlcolor   = linkblue,
]{hyperref}
\onehalfspacing

% S-prefixed figures, tables, equations
\renewcommand{\thefigure}{S\arabic{figure}}
\renewcommand{\thetable}{S\arabic{table}}
\renewcommand{\theequation}{S\arabic{equation}}
\renewcommand{\thesection}{S\arabic{section}}

% Running header
\pagestyle{fancy}
\fancyhf{}
\renewcommand{\headrulewidth}{0.4pt}
\fancyhead[L]{\small\textit{Supplementary Material}}
\fancyhead[R]{\small\textit{Huff Curves for Brazil}}
\fancyfoot[C]{\small\thepage}
\fancypagestyle{plain}{%
  \fancyhf{}
  \renewcommand{\headrulewidth}{0pt}
  \fancyfoot[C]{\small\thepage}
}

\title{\textbf{Supplementary Material}\\[6pt]
{\large Empirical Huff Curves for Brazil: A National-Scale Derivation
from ANA Sub-Daily Rainfall Records}}

\author{Marcus Nobrega Gomes Junior}
\date{}

\begin{document}
\maketitle

\begin{center}\small
  \textit{$^1$Department of Earth System Science, Stanford University,
  Braun Hall, Bldg.\ 320, 450 Jane Stanford Way,
  Stanford, California 94305, USA}\\[3pt]
  \textit{$^2$Department of Hydrology and Atmospheric Sciences,
  University of Arizona,
  John W.\ Harshbarger Building, Tucson, Arizona 85719, USA}\\[5pt]
  ORCID: 0000-0002-8250-8195\quad
  \texttt{mngomes@stanford.edu}
\end{center}
\vspace{0.5em}
\noindent\rule{\linewidth}{0.4pt}
\vspace{0.5em}

This document provides general supporting information for the paper.
Section~\ref{sec:supp_pipeline} summarises the station processing
pipeline and quality-control outcome for the full 3{,}164-station
ANA network.
Sections~\ref{sec:supp_purpose}--\ref{sec:supp_results} document the
design-hydrograph sensitivity experiment in full detail, covering
the experimental design (Section~\ref{sec:supp_purpose}), input
datasets (Section~\ref{sec:supp_data}), catchment selection
(Section~\ref{sec:supp_catch}), and the complete mathematical
description of the rainfall--runoff method
(Section~\ref{sec:supp_math}), so that the experiment can be
reproduced independently.
Section~\ref{sec:supp_quartile} provides the per-station quartile
distribution for all accepted stations.

\vspace{1em}
\noindent\rule{\linewidth}{0.4pt}
{\setlength{\parskip}{0pt}
\listoftables
\listoffigures}
\noindent\rule{\linewidth}{0.4pt}
\vspace{0.5em}

%%====================================================================
\section{Processing pipeline and quality-control summary}
\label{sec:supp_pipeline}

Table~\ref{tab:supp_processing} summarises the outcome of the
station-level quality-control pipeline applied to the 3{,}164 ANA
telemetric stations. A station was accepted if it met all four
criteria simultaneously: (i) record span $\geq$~4~yr; (ii) missing
fraction $\leq$~20\%; (iii) no calendar year with a full-zero daily
record (instrument failure indicator); and (iv) no interval depth
implying an intensity exceeding 300~mm~h$^{-1}$.

\begin{table}[h]
  \centering
  \caption{Processing pipeline outcome for the 3{,}164 ANA telemetric
    stations. Quality filters applied simultaneously: record span
    $\geq$~4~yr; missing fraction $\leq$~20\%; no full-zero calendar
    year; maximum interval intensity $\leq$~300~mm~h$^{-1}$.}
  \label{tab:supp_processing}
  \small
  \begin{tabular}{lrrp{6.5cm}}
    \toprule
    Outcome & $N$ & \% & Notes \\
    \midrule
    Accepted — curves fitted & 1{,}045 & 33.0 &
      Passed all filters; empirical Huff curves derived \\
    Rejected — quality       & 2{,}086 & 65.9 &
      Failed $\geq$1 filter (93.1\% excessive missing data;
      1.9\% full-zero calendar year; 1.5\% insufficient
      record length) \\
    Rejected — no data       &    33   &  1.0 &
      No locally cached ANA data available \\
    \midrule
    \textbf{Total}           & \textbf{3{,}164} & \textbf{100.0} & \\
    \bottomrule
  \end{tabular}
\end{table}

%%====================================================================
\section{Purpose and design of the experiment}
\label{sec:supp_purpose}

The experiment is a \emph{controlled} comparison: for each catchment a
single design hydrograph is computed twice, once with the original
Huff (1967) first-quartile (Q1) reference curve and once with the
locally derived biome Q1 curve from this study. The catchment geometry,
curve number, time of concentration, design rainfall depth, and storm
duration are held identical between the two runs. The only quantity
that differs is the dimensionless temporal pattern (the cumulative mass
curve), so the resulting change in peak discharge $\Delta Q_p$ and in
time-to-peak $\Delta t_p$ is attributable solely to the hyetograph
shape. Because the loss model and unit hydrograph are common to both
runs, systematic errors in the absolute curve number, time of
concentration, and unit-hydrograph shape largely cancel in the relative
comparison.

%%====================================================================
\section{Input datasets}
\label{sec:supp_data}

Table~\ref{tab:supp_data} lists the datasets used. Their spatial fields
are shown in Fig.~\ref{fig:supp_inputs}: the curve number, time of
concentration, and 25-year design depth at the sampled catchments
(panels a--c), the soil clay fraction (panels d--e), and the gridded
1-hour, 25-year design intensity derived from the gridded Sherman IDF
surfaces of \citet{GomesJuniorIDF} (panel f).

\begin{table}[h]
  \centering
  \caption{Input datasets for the design-hydrograph experiment.}
  \label{tab:supp_data}
  \small
  \begin{tabular}{p{2.6cm}p{3.2cm}p{2.0cm}p{4.4cm}}
    \toprule
    Quantity & Source & Resolution & Use \\
    \midrule
    Catchments & HydroBASINS level 12 \citep{Lehner2013} & vector &
      Catchment polygons, drainage area \\
    Terrain & Copernicus GLO-90 DEM \citep{CopernicusDEM} & 90~m &
      Relief and channel slope $\rightarrow$ time of concentration \\
    Soil texture & SoilGrids 2.0 \citep{Poggio2021} & 250~m &
      Clay/sand fractions $\rightarrow$ hydrologic soil group $\rightarrow$ CN \\
    Design rainfall & Gridded Sherman IDF \citep{GomesJuniorIDF} &
      0.1$^{\circ}$ & Design depth at duration $=T_c$, $T=25$~yr \\
    Temporal pattern & This study (biome Q1) and Huff (1967) Q1 &
      --- & Dimensionless cumulative mass curves \\
    \bottomrule
  \end{tabular}
\end{table}

\begin{figure}[h]
  \centering
  \includegraphics[width=\textwidth]{figures/supp_inputs}
  \caption{Input-data fields for the 579 catchments. (a) Curve number
    (antecedent moisture condition II); (b) time of concentration;
    (c) 25-year design rainfall depth for a storm of duration equal to
    the time of concentration; (d) soil clay fraction at the catchments;
    (e) gridded SoilGrids 250~m clay fraction; (f) gridded 1-hour,
    25-year design rainfall intensity computed from the Sherman IDF
    parameters. Biome (light grey) and national (dark grey) boundaries
    are shown.}
  \label{fig:supp_inputs}
\end{figure}

%%====================================================================
\section{Catchment selection}
\label{sec:supp_catch}

The SCS-CN method and the SCS dimensionless unit hydrograph were
developed for small, infiltration-excess (Hortonian) agricultural
watersheds and are not valid where saturation-excess (variable
source area) runoff dominates or where rainfall cannot be treated as
spatially uniform \citep{NRCS2004,TR55_1986}. Catchments were therefore
screened from the HydroBASINS level-12 layer using the criteria in
Table~\ref{tab:supp_criteria}. Of 30{,}866 eligible catchments, a
stratified random sample of 150 per biome (600 total) was drawn;
579 yielded complete input data and a valid hydrograph
(Fig.~\ref{fig:supp_catchments}).

\begin{table}[h]
  \centering
  \caption{Catchment eligibility criteria.}
  \label{tab:supp_criteria}
  \small
  \begin{tabular}{p{4.2cm}p{8.0cm}}
    \toprule
    Criterion & Rationale \\
    \midrule
    Drainage area 5--250~km$^2$ &
      Small enough for a lumped unit hydrograph and spatially uniform
      design rainfall \\
    Biome $\in$ \{Cerrado, Caatinga, Mata Atl\^antica, Pampa\} &
      Infiltration-excess regimes; Amaz\^onia and Pantanal excluded as
      saturation-excess / wetland \\
    Valid terrain, soil, and IDF data &
      Finite relief, soil texture, and Sherman parameters at the
      catchment \\
    Time of concentration $0.1$--$24$~h &
      Within the validity window of the IDF and unit-hydrograph methods \\
    \bottomrule
  \end{tabular}
\end{table}

\begin{figure}[h]
  \centering
  \includegraphics[width=0.66\textwidth]{figures/supp_catchments}
  \caption{Eligible HydroBASINS level-12 catchments (coloured by biome)
    and the 579 sampled catchments used in the experiment (black points).}
  \label{fig:supp_catchments}
\end{figure}

The curve number (antecedent moisture condition II) was assigned from
the catchment hydrologic soil group (HSG) and the biome-characteristic
land cover, using the values in Table~\ref{tab:supp_cn}. The HSG was
derived from the SoilGrids clay and sand fractions $f_{\mathrm{clay}}$
and $f_{\mathrm{sand}}$ (\%): group A if $f_{\mathrm{sand}}\!\ge\!70$
and $f_{\mathrm{clay}}\!<\!10$; D if $f_{\mathrm{clay}}\!\ge\!40$;
C if $f_{\mathrm{clay}}\!\ge\!27$; B if $f_{\mathrm{clay}}\!<\!18$ and
$f_{\mathrm{sand}}\!\ge\!50$; otherwise C if $f_{\mathrm{clay}}\!\ge\!20$
and B if not.

\begin{table}[h]
  \centering
  \caption{Curve number (AMC II) by biome land cover and hydrologic
    soil group \citep{NRCS2004}.}
  \label{tab:supp_cn}
  \small
  \begin{tabular}{lcccc}
    \toprule
    Biome (land-cover archetype) & A & B & C & D \\
    \midrule
    Cerrado (savanna/woodland, good)      & 35 & 58 & 70 & 77 \\
    Caatinga (semi-arid shrub, fair)      & 45 & 66 & 77 & 83 \\
    Mata Atl\^antica (forest/ag mosaic)   & 36 & 60 & 73 & 79 \\
    Pampa (grassland/pasture, good)       & 39 & 61 & 74 & 80 \\
    \bottomrule
  \end{tabular}
\end{table}

%%====================================================================
\section{Mathematical description of the method}
\label{sec:supp_math}

\subsection{Design rainfall depth (Sherman IDF)}

The design rainfall intensity $i$ (\si{mm/h}) is obtained from the
gridded Sherman intensity--duration--frequency relationship of
\citet{GomesJuniorIDF},
\begin{equation}
  i = \frac{K\,T^{a}}{(t + b)^{c}},
  \label{eq:sherman}
\end{equation}
where $T$ is the return period (yr; here $T=25$), $t$ is the storm
duration (min), and $K$, $a$, $b$, $c$ are the gridded Sherman
parameters sampled at the catchment centroid. The storm duration is set
equal to the time of concentration, $t = 60\,T_c$ (with $T_c$ in
hours), and the total design depth is
\begin{equation}
  P = i\,(t = 60\,T_c)\;\times\;T_c .
  \label{eq:depth}
\end{equation}

\subsection{Time of concentration}

The main-channel length is estimated from drainage area $A$
(\si{km^2}) by Hack's law,
\begin{equation}
  L = 1.4\,A^{0.6}\quad[\si{km}],
  \label{eq:hack}
\end{equation}
and the channel slope is taken as the catchment relief (5th-to-95th
percentile elevation range $\Delta z$, from the Copernicus DEM) over the
channel length, $S = \Delta z / L$. The time of concentration follows
\citet{Kirpich1940},
\begin{equation}
  T_c = \frac{1}{60}\,\bigl[\,0.0078\,L_{\mathrm{ft}}^{0.77}\,S^{-0.385}\,\bigr]
  \quad[\si{h}],
  \label{eq:kirpich}
\end{equation}
with $L_{\mathrm{ft}} = L/0.3048$ the channel length in feet, and
$T_c$ clipped to the range $0.1$--$24$~h.

\subsection{Design hyetograph from the Huff curve}

The storm is discretised into $n$ steps of length $\Delta t$. With
$\tau_k = t_k/(60\,T_c)$ the normalised time and $F(\tau)$ the
dimensionless cumulative Huff mass curve (the 7th-degree polynomial,
clipped to $[0,1]$, monotonised, with endpoints fixed to $\{0,1\}$),
the incremental rainfall depth in step $k$ is
\begin{equation}
  \Delta P_k = \bigl[\,F(\tau_k) - F(\tau_{k-1})\,\bigr]\,P .
  \label{eq:hyeto}
\end{equation}
This is evaluated separately for the original Huff (1967) Q1 curve and
for the biome Q1 curve of this study.

\subsection{SCS-CN direct runoff}

The maximum potential retention $S$ (\si{mm}) is obtained from the
curve number,
\begin{equation}
  S = \frac{25400}{\mathrm{CN}} - 254 ,
  \label{eq:retention}
\end{equation}
and the initial abstraction is $I_a = \lambda S$ with $\lambda = 0.2$.
The cumulative direct runoff at cumulative rainfall $P_{\mathrm{cum}}$
is
\begin{equation}
  Q_{\mathrm{cum}} =
  \begin{cases}
    \dfrac{(P_{\mathrm{cum}} - I_a)^2}{P_{\mathrm{cum}} - I_a + S},
      & P_{\mathrm{cum}} > I_a,\\[8pt]
    0, & \text{otherwise,}
  \end{cases}
  \label{eq:scscn}
\end{equation}
and the incremental excess (effective) rainfall in step $k$ is the first
difference $\Delta Q_k = Q_{\mathrm{cum},k} - Q_{\mathrm{cum},k-1}$.

\subsection{SCS dimensionless unit hydrograph}

The basin lag is $L_g = 0.6\,T_c$ and, for a unit-rainfall duration
equal to the computational step $D=\Delta t$, the time to peak of the
unit hydrograph is
\begin{equation}
  T_p = \frac{D}{2} + L_g .
  \label{eq:tp}
\end{equation}
The unit-hydrograph peak discharge (per mm of excess rainfall) uses the
SI form of the standard peak-rate factor (PRF~484),
\begin{equation}
  q_p = \frac{0.208\,A}{T_p}\quad[\si{m^3.s^{-1}.mm^{-1}}],
  \label{eq:qp}
\end{equation}
and the unit-hydrograph ordinates are $u(t) = q_p\,g(t/T_p)$, where
$g(\cdot)$ is the tabulated SCS dimensionless unit hydrograph
(PRF~484). The outlet hydrograph is the convolution of the incremental
excess hyetograph with the unit hydrograph,
\begin{equation}
  Q(t_j) = \sum_{k} \Delta Q_k \, u(t_j - t_k).
  \label{eq:convol}
\end{equation}

\subsection{Comparison metrics}

For each catchment, Eqs.~\eqref{eq:sherman}--\eqref{eq:convol} are
evaluated with the original Huff (1967) Q1 curve (reference, subscript
$r$) and with the biome Q1 curve (local, subscript $l$). The reported
metrics are the relative change in peak discharge and the shift in
time-to-peak,
\begin{equation}
  \Delta Q_p = 100\,\frac{Q_{p,l} - Q_{p,r}}{Q_{p,r}}\;[\%],
  \qquad
  \Delta t_p = t_{p,l} - t_{p,r}\;[\si{h}].
  \label{eq:metrics}
\end{equation}
Positive $\Delta Q_p$ indicates that the locally derived curve produces
a higher design peak than the imported reference; negative $\Delta t_p$
indicates an earlier peak.

%%====================================================================
\section{Per-catchment results}
\label{sec:supp_results}

The complete per-catchment table — HydroBASINS identifier, biome,
coordinates, drainage area, relief, slope, time of concentration, soil
texture, hydrologic soil group, curve number, design depth, reference
and local peak discharges, $\Delta Q_p$, and $\Delta t_p$ — and the
design hydrographs are provided in the data archive
(\texttt{outputs/hydro\_impact/}).

%%====================================================================
\section{Per-station quartile distribution}
\label{sec:supp_quartile}

The complete event-quartile composition of every accepted station is
provided as a data file,
\texttt{station\_quartile\_distribution.csv} (1{,}045 rows), with the
columns: station identifier, latitude, longitude, state, biome, record
length, total event count, the event count and percentage in each
quartile (Q1--Q4), and the dominant quartile. Table~\ref{tab:supp_quartile}
shows one representative station per biome (the station with the longest
event record in each) to illustrate the format and the range of
compositions; the full distribution is summarised at biome level in
Table~2 of the main text and mapped in Fig.~6.

\begin{table}[h]
  \centering
  \caption{Event-quartile distribution for the longest-record station in
    each biome (illustrative excerpt of the full per-station data file).
    Q1--Q4 give the percentage of the station's events in each quartile;
    the final column is the dominant (modal) quartile.}
  \label{tab:supp_quartile}
  \small
  \setlength{\tabcolsep}{5pt}
  \begin{tabular}{lllrrrrrc}
    \toprule
    Station & State & Biome & $N$ events &
      Q1 (\%) & Q2 (\%) & Q3 (\%) & Q4 (\%) & Dom. \\
    \midrule
    14100000 & AM & Amaz\^onia       & 777 & 55.0 & 16.2 & 15.4 & 13.4 & 1 \\
    39540550 & PE & Caatinga         & 270 & 29.6 & 29.3 & 21.5 & 19.6 & 1 \\
    27110000 & TO & Cerrado          & 511 & 55.6 & 19.2 & 13.9 & 11.4 & 1 \\
    80360200 & SP & Mata Atl\^antica & 666 & 23.7 & 22.4 & 17.9 & 36.0 & 4 \\
    74120700 & RS & Pampa            & 407 & 36.4 & 31.9 & 22.1 &  9.6 & 1 \\
    66910000 & MS & Pantanal         & 385 & 46.0 & 23.6 & 17.9 & 12.5 & 1 \\
    \bottomrule
  \end{tabular}
\end{table}

\bibliographystyle{elsarticle-harv}
\bibliography{references}

\end{document}
